%! AP Statistics — Unit 1, Lesson 1.5 — Exit Ticket ANSWER KEY
\documentclass[10pt]{article}
\usepackage{apstats-article}
\usepackage{apstats-key}

%=============================================================================
\begin{document}

\pageheader{Unit 1, Lesson 1.5}{Exit Ticket --- Answer Key}
\begin{tabularx}{\linewidth}{X X X}
Name:\;\ans{Key} & Date:\; & Period:\;
\end{tabularx}

\vspace{0.18in}

\begin{tcolorbox}[colback=sky, colframe=navy, boxrule=0.9pt, arc=2mm,
    left=4mm, right=4mm, top=3mm, bottom=3mm]
\small A nurse recorded the resting pulse rates (beats per minute) of
\textbf{12 patients}. The stem-and-leaf plot below displays the data.

\vspace{0.1in}
\begin{center}
\renewcommand{\arraystretch}{1.3}
\begin{tabular}{r | l}
\textbf{Stem} & \textbf{Leaf} \\
\hline
6 & 2\; 4\; 5\; 8 \\
7 & 0\; 2\; 2\; 4\; 5\; 6\; 8 \\
8 & 0 \\
\end{tabular}
\quad Key: \textbf{7 $|$ 2} = 72 bpm
\end{center}
\end{tcolorbox}

\vspace{0.2in}

\begin{enumerate}[leftmargin=1.5em, itemsep=10pt]

  \item \ans{7 patients} have a pulse rate in the 70s.

  \begin{teachernote}
  Count leaves in the 7-stem row: 0,2,2,4,5,6,8 — that is 7 values.
  \end{teachernote}

  \item \ansline{Range $= 80 - 62 = \mathbf{18}$ bpm.}

  \item \ansline{$n = 12$, so the median is the average of the 6th and 7th values
        in order. Sorted: 62, 64, 65, 68, 70, 72, 72, 74, 75, 76, 78, 80.
        Median $= (72 + 72)/2 = \mathbf{72}$ bpm.}

  \begin{teachernote}
  Some students may count 6 and 7 from the stemplot directly without listing
  all values. Both approaches are valid; require that the position (6th and 7th)
  is identified.
  \end{teachernote}

  \item \ans{Skewed left} --- the distribution has a longer tail toward the lower values
        (60s) relative to the cluster in the 70s and 80s.

  \begin{teachernote}
  The single 80 and the cluster in the 70s, with the spread-out 60s values, create
  a leftward tail. Some students may argue approximately symmetric; this is marginally
  defensible but skewed left is the better answer with $n=12$.
  \end{teachernote}

\end{enumerate}

\vspace{0.14in}

\begin{tcolorbox}[colback=redbg, colframe=redacc, boxrule=0.8pt, arc=2mm,
    left=4mm, right=4mm, top=3mm, bottom=3mm]
\small\textbf{AP-Style Practice --- Key}\\[8pt]
\ans{The distribution of resting pulse rates for these 12 patients is skewed left,
with a typical value around 72 bpm (the median) and a range of 18 bpm from 62 to
80 bpm.}

\vspace{8pt}
\begin{teachernote}
Require: (1) a shape description (skewed left or approximately symmetric with
justification); (2) a center value stated in context with units; (3) a spread
measure (range, or min to max) with units. Deduct for missing units (bpm) or
missing context.
\end{teachernote}
\end{tcolorbox}

\end{document}
